%==============================================================================
\section{Discussion}
\label{sec:discussion}

\subsection{Which architectural family wins on CAMUS?}
\label{sec:disc-families}

The top tier of our benchmark is mixed across paradigms, and no single
architecture wins on every axis. The hybrid CNN--Transformer TransUNet leads on
bulk Dice and HD95 and holds the highest EF correlation; the pure CNN UNet-V1
attains the lowest EF MAE ($6.62\%$) and the tightest Bland--Altman bias
($-0.42\%$); and nnU-Net leads on parameter efficiency while sitting mid-table
on EF at $8.32\%$. The difference between TransUNet and nnU-Net on mean Dice is
$0.0023$, below the Bonferroni-corrected significance threshold.

A practitioner choosing among them should therefore be guided not by Dice but
by the axis they actually care about: deployment constraints favour nnU-Net at
$6.5\times$ fewer parameters, clinical EF agreement favours UNet-V1, and
contour precision favours TransUNet. That these three come apart is itself the
finding --- a single leaderboard column would have concealed it.

The pure-Transformer family (Swin-UNet) is competitive on bulk Dice
($0.8882$) but visibly worse on EF (MAE $9.44\%$, $r = 0.640$ --- the lowest
correlation of any architecture in the benchmark). This
finding is consistent with reports that pure-attention models, despite
strong overall overlap, can produce boundary contours that are subtly
biased \citep{chen2021transunet}, an effect that compounds when
contours feed Simpson's biplane volumetric integration.

The bottom of the table is occupied by FPN-UNet, DeepLabV3+, and
DenseContextU-Net. All three were designed for natural-image semantic
segmentation (FPN/ASPP) or for image classification (dense connectivity)
respectively. Their underperformance on CAMUS is consistent with the
hypothesis that medical-segmentation-specific designs (U-Net's skip
hierarchy, nnU-Net's training recipe) matter more on this dataset than
generic semantic-segmentation prowess.

\subsection{Comparison to the published CAMUS baselines}
\label{sec:disc-leaderboard}

The original CAMUS paper \citep{leclerc2019} reported nnU-Net HD95
$\approx 4.3$~mm and EF MAE $\sim$$9\%$. Our reproduced nnU-Net under the
unified protocol reaches HD95 $4.27$~mm and EF MAE $8.32\%$
($r = 0.783$).
Both axes
therefore behave the same way, and it is not the way we first read them:
boundary accuracy is statistically indistinguishable from the 2019 result,
and so, very nearly, is EF.

This is a negative result about protocol, and we state it as such because
an earlier version of this analysis reported the opposite. Computing EF from
predictions on the resized $256 \times 256$ grid while passing the native
pixel spacing yields ventricular volumes near $10$~ml against a true
$\approx 108$~ml, and the error does not cancel in the EF ratio. Every EF
number in this paper is now computed at native resolution with the official
CAMUS implementation; Section~\ref{sec:results-clinical} reports the result, and the
ground-truth oracle at $5.44\%$ bounds it.

We therefore read the contour \emph{and} the derived clinical quantity as
close to their respective ceilings on CAMUS --- the first set by annotation
variability, which is of the same order as the differences we measure between
architectures, and the second by the biplane estimator itself. Three
training-protocol upgrades made standard since 2019 are nonetheless worth
recording, because they are what makes the reproduction possible at all:
\begin{enumerate}[leftmargin=*,nosep]
\item Cosine-with-warm-restarts learning-rate schedule with
      $10$-epoch warmup, in place of the original ReduceLROnPlateau.
\item Mixed-precision training (FP16 AMP), which allows larger effective
      batch sizes and smoother loss landscapes.
\item Composite loss with an explicit boundary term
      \citep{kervadec2019}, which directly penalises the
      Hausdorff-relevant boundary error.
\end{enumerate}
This is a useful baseline calibration point for the field: cross-paper
comparisons of CAMUS results from before $\sim 2021$ to after should be
read with the understanding that protocol changes alone produce
sub-millimetre HD95 improvements at constant architecture.

\subsection{A boundary-metric pitfall, and what it implies for the
literature}
\label{sec:disc-pitfall}

The correction described in Section~\ref{sec:results-boundary} is worth
stating as a general point, because nothing about it is specific to our
models. Segmentation networks consume a fixed input size; clinical images
do not have one. Any pipeline that resizes an image for inference, scores
the prediction on that resized grid, and then multiplies pixel distances
by the spacing recorded in the original DICOM or NIfTI header will
under-report every boundary distance by the resampling factor. Overlap
metrics are immune --- Dice and IoU are ratios and cancel the scale ---
which is precisely why the error survives casual inspection: a pipeline
can produce entirely correct Dice and badly wrong HD95 from the same
predictions, with no internal inconsistency to notice.

On CAMUS the factor is $\approx 2.2$ (native acquisitions average
$630 \times 519$ against a $256 \times 256$ network grid), and the
consequence is not cosmetic. It moves our best model from $1.82$~mm to
$4.08$~mm, and thereby converts ``roughly twice as accurate as the
published baseline'' into ``indistinguishable from the published
baseline.'' Both numbers come from the same weights and the same test
split; only the measurement convention differs.

We therefore suggest reading sub-$2$~mm HD95 values on CAMUS with
caution, and we make no claim about which published results are affected
--- establishing that would require the evaluation code for each, which is
usually unavailable. The constructive recommendation is narrower and
checkable: report the resolution at which boundary metrics were computed,
and prefer acquisition resolution, since that is the only convention under
which a millimetre means the same thing across papers. This is a concrete
instance of the reporting problems documented by
\citet{reinke2021limitations} and \citet{maierhein2024metrics}, and it is
consistent with the broader finding of \citet{isensee2024revisited} and
\citet{kazaj2025claims} that apparent architectural progress in medical
segmentation frequently does not survive standardised validation. Our own
benchmark is an instance of that pattern rather than an exception to it:
once measurement is held fixed, nine architectures spanning three
paradigms separate by less than the noise floor on the boundary axis.

\subsection{Design choices that matter}
\label{sec:disc-design}

Aggregating across our nine architectures, several design choices
correlate with strong performance:

\begin{itemize}[leftmargin=*,nosep]
\item \textbf{Pretrained encoders.} Five of the nine architectures
      initialise their encoder from ImageNet or BiT weights (TransUNet,
      UNet-ResNet, Swin-UNet, FPN-UNet, DeepLabV3+; Table~\ref{tab:archs}).
      Pretraining is nonetheless neither necessary nor sufficient for
      top-tier accuracy on CAMUS: the two scratch-trained models nnU-Net and
      UNet-V1 place second and fourth on mean Dice --- ahead of the
      pretrained Swin-UNet, FPN-UNet, and DeepLabV3+ --- and UNet-V1 also
      attains the best EF agreement of any architecture tested. Encoder
      pretraining thus appears to help mainly where
      an architecture's from-scratch inductive bias is weaker, rather than
      conferring a uniform advantage on this dataset.
\item \textbf{Deep skip connections} (every U-Net-shaped architecture
      except DeepLabV3+). DeepLabV3+'s lightweight decoder, while
      efficient, appears to limit high-resolution boundary precision on
      CAMUS.
\item \textbf{Boundary-aware loss} (used uniformly in our protocol).
      This is a protocol choice rather than an architecture choice. Since
      our HD95 matches rather than beats the published baseline, we make
      no claim that it improves absolute boundary accuracy; its role here
      is to hold the boundary objective constant across all nine
      architectures so that the comparison is like-for-like.
\item \textbf{Auto-configuration} (nnU-Net's distinguishing feature).
      The nnU-Net framework produces a competitive model on this dataset
      with $\sim$$15$~M parameters --- substantially smaller than any
      pretrained-encoder architecture --- by selecting patch size, batch
      size, and augmentation strategy from data statistics.
\end{itemize}

\subsection{Clinical implications}
\label{sec:disc-clinical}

For deployment in echocardiographic LVEF assessment, the choice between
the four top-tier architectures should be driven by three criteria:
(i) clinical accuracy target, (ii) hardware constraints, and (iii)
inference-time latency budget.
\begin{itemize}[leftmargin=*,nosep]
\item For \emph{tightest EF agreement}, UNet-V1 attains the lowest EF MAE
      ($6.56\%$) and TransUNet the smallest Bland--Altman bias
      ($-0.54\%$). The two differ by less than the run-to-run variation
      documented in Sec.~\ref{sec:results-clinical}, so we do not
      recommend one over the other on this axis.
\item For \emph{maximum boundary accuracy} at deployment-irrelevant cost,
      TransUNet has the lowest HD95 ($4.08$~mm).
\item For \emph{edge / portable ultrasound deployment} where
      parameters and latency dominate, nnU-Net at $15.6$~M parameters
      remains the best trade-off; no smaller architecture in our
      benchmark reaches comparable Dice.
\end{itemize}
None of the four top-tier models reach below-$5\%$ EF MAE in this study,
indicating that human-level cardiologist agreement is still a target
boundary for the field rather than a solved problem.

\subsection{Limitations and future work}
\label{sec:disc-limits}

We acknowledge five limitations.

\begin{enumerate}[leftmargin=*,nosep]
\item \textbf{Single-dataset.} All evaluation is on CAMUS. A cross-dataset
      transfer evaluation (\emph{e.g.,} EchoNet-Dynamic
      \citep{ouyang2020}) would strengthen the generality claims and is
      planned for future work.
\item \textbf{2D only.} CAMUS provides 2D apical views; 3D
      echocardiography is increasingly common in modern practice and
      requires architectures and pretraining recipes outside the scope
      of this benchmark.
\item \textbf{Test set size.} The CAMUS official test set is 50 patients
      ($200$ ED+ES frames). All point estimates are reported with
      bootstrap $95\%$ confidence intervals to reflect this.
\item \textbf{No SSM models.} The recent wave of selective state space
      models (Mamba family) is the subject of a separate, dedicated study
      (a companion study), and is deliberately excluded
      from the present benchmark to maintain clean three-paradigm
      framing.
\item \textbf{Single training protocol.} A single shared protocol is the
      goal of this benchmark, but some architectures may benefit from
      protocol-specific tuning (\emph{e.g.,} larger batch sizes for
      pretrained Transformers). Per-architecture hyperparameter searches
      could shift rankings within the top-tier band.
\end{enumerate}

\subsection{Take-home messages}
\label{sec:disc-takehome}

\begin{enumerate}[leftmargin=*,nosep]
\item No single architectural family dominates CAMUS. Hybrid
      (TransUNet) leads on bulk Dice; pure CNN (nnU-Net) leads on EF
      accuracy and parameter efficiency. Pure Transformer (Swin-UNet) is
      competitive on Dice but degraded on EF.
\item Training-protocol modernisation (cosine schedule, AMP, boundary
      loss) accounts for most of the $\sim$$2\times$ improvement in
      HD95 over the original CAMUS baselines, independent of
      architecture.
\item LV-epicardium remains the hardest of the three foreground classes
      across every architectural family.
\item The accuracy--parameter Pareto frontier is dominated by nnU-Net at
      the efficient end and TransUNet at the high-accuracy end.
\end{enumerate}
